\begin{theorem}[A composed two-step physical endpoint route]
    \label{thm:peps_two_step_string_physical_route}
    \lean{TNLean.PEPS.torusSweptStringRightStep_extension_initial,TNLean.PEPS.IsGIsometric.exists_unitary_torusTwoStepStringRoute}
    \leanok
    Let the native torus periods satisfy $w\ge8$, $H\ge7$, and let the
    original site tensor be regular $G$-isometric. Denote the actual closed
    states of the initial, rightward-moved, and subsequently upward-moved
    assignments by $\psi_0(g,h),\psi_1(g,h),\psi_2(g,h)$.
    There exist original-spin unitaries $W_1,W_2$ on the respective six-site
    tiles, chosen before every $g,h\in G$, whose extensions $T_1,T_2$ by the
    identity satisfy
    \begin{align}
        T_1\psi_0(g,h)&=\psi_1(g,h), &
        T_2\psi_1(g,h)&=\psi_2(g,h), &
        (T_2T_1)\psi_0(g,h)&=\psi_2(g,h).
    \end{align}
    Each extension and their product are unitary. The intermediate bond
    assignment is the literal one-bond update of the initial assignment, with
    the common exterior retained. This is an auxiliary two-step part of
    \cite[Section~6.6]{Schuch2010PEPS}; it does not assert the full braid or
    the partner-reunion operation.
\end{theorem}
\begin{proof}
    \leanok
    \uses{thm:peps_initial_string_right_physical_step,thm:peps_second_string_physical_crossing}
    The bond changed by the rightward step is internal to its actual tile, so
    extending the rightward assignment by the initial exterior gives exactly
    the literal intermediate assignment. Select the two regional unitaries
    before the labels and apply their closed-state identities successively.
    The product remains unitary.
\end{proof}
